\documentclass[10pt,a4paper]{amsart}
\usepackage[margin=19mm]{geometry}
\usepackage{amsmath,amssymb,mathtools}
\usepackage[scr=rsfs,cal=euler]{mathalfa}
\usepackage{xfrac}
\usepackage[unicode,hidelinks,bookmarks=false]{hyperref}
\newcommand{\ie}{i.e.\ }
\newcommand{\cf}{cf.\ }
\newcommand{\resp}{respectively\ }
\newcommand{\Subsection}{\subsection}
\providecommand{\nobreakdash}{\nobreak\hbox{-}}
\setlength{\emergencystretch}{2em}
\multlinegap0pt
\usepackage{bm}
\usepackage{bbm}
\usepackage{dsfont}
\usepackage{xspace}
\usepackage{cancel}
\usepackage{mathtools}
\usepackage{amsthm}
\makeatletter
\@ifundefined{theorem}{\newtheorem{theorem}{Theorem}}{}
\@ifundefined{corollary}{\newtheorem{corollary}[theorem]{Corollary}}{}
\@ifundefined{example}{\newtheorem{example}[theorem]{Example}}{}
\makeatother
\newcommand{\SM}{{\setminus}}
\newcommand{\w}{\omega}
\title[New classes of compact-type spaces]{New classes of compact-type spaces}
\author{Saak Gabriyelyan, Evgenii Reznichenko}
\date{Source: arXiv:2510.21642v1 (CC BY 4.0)}
\begin{document}
\maketitle

\noindent\emph{Source and licence.} New classes of compact-type spaces. Version arXiv:2510.21642v1 (CC BY 4.0), distributed under the Creative Commons Attribution 4.0 International licence, as recorded in the arXiv licence metadata for this exact version and shown on the abstract page https://arxiv.org/abs/2510.21642v1. Full rights notice: licenses/arxiv-2510.21642-CC-BY-4.0.txt.\par\smallskip

\noindent\textit{Editorial scope.} Counted unit: the Example whose source label is \texttt{exa:Veer-Q} in the article (source0.tex lines 1140--1142, source label \texttt{exa:Veer-Q}), delivered together with the article's own complete original proof of it (source0.tex lines 1145--1180, one \texttt{proof} environment of 36 lines). No mathematics is written, repaired or re-derived: statement and proof are the article's own text line for line. Required context retained uncounted: c:counable-X-woca-Seq-Ascoli (lines 779--785). Every definition, display and label of those lines is retained unchanged. Omitted from this package and not redistributed: 1--778, 786--1139, 1143--1144, 1181--1880. Nothing inside a retained excerpt is omitted or altered: every retained body is delivered exactly as the article's lines are written, so no omission receipt of any kind accompanies this package. Citations: the retained excerpts carry no citation command at all, so no bibliography entry of the article is transcribed and no reference is redistributed. Reference adaptation: none was needed and none was made; every reference inside the delivered unit resolves inside it, so no unresolved reference or citation is produced. Printed slips kept literal: no printed word, symbol, hypothesis or number of the counted statement or of its proof was altered. Layout and class: the article's own class is \texttt{elsarticle} with packages including amsfonts, amsmath, amscd, amssymb, pb-diagram, xy, color; the wrapper is the lane's pinned \texttt{amsart} editorial wrapper at 10pt on A4, which restates the theorem-like environments the retained text needs and adds the article's own macro definitions that the retained text needs (\texttt{\textbackslash SM}, \texttt{\textbackslash w}), with extra wrapper packages (bm, bbm, dsfont, xspace, cancel, mathtools). The delivered numbering is the unit's own. Assets: no retained span contains a figure environment, a graphics-inclusion command, a PDF-page inclusion, a table or a TikZ picture; no image file is copied into this package, the package declares no asset dependency and no image file is redistributed with it. Licence: version arXiv:2510.21642v1 of this article is distributed under the Creative Commons Attribution 4.0 International licence, as recorded in the arXiv licence metadata for this exact version and shown on the abstract page https://arxiv.org/abs/2510.21642v1. A full rights notice is delivered at licenses/arxiv-2510.21642-CC-BY-4.0.txt. Characters: every retained excerpt is pure ASCII, so the delivered engine, XeLaTeX with its default Latin Modern text font, reports no missing character.
\begin{corollary} \label{c:counable-X-woca-Seq-Ascoli}
Let $X$ be a countable space. Then:
\begin{enumerate}
\item[{\rm(i)}] $X$ is $\kappa$-sequential if, and only if, $X$ is weakly open-compact attainable if, and only if,  $X$ is a sequentially Ascoli space;
\item[{\rm(ii)}] $X$ is $\kappa$-Fr\'{e}chet--Urysohn if, and only if, it is open-compact attainable.
\end{enumerate}
\end{corollary}

\begin{example} \label{exa:Veer-Q}
The product $V(\w)\times \mathbb{Q}$ is neither weakly open-compact attainable nor sequentially Ascoli.
\end{example}

\begin{proof}
Since $X:=V(\w)\times \mathbb{Q}$ is countable, by (i) of Corollary \ref{c:counable-X-woca-Seq-Ascoli}, it suffices to prove that $X$ is not $\kappa$-sequential.

Choose two sequences $\{a_k\}_{k\in\w}$ and  $\{b_k\}_{k\in\w}$ of irrational numbers such that
\begin{enumerate}
\item[(a)] $0<b_{k+1}<a_k<b_k$ for every $k\in\w$;
\item[(b)] $b_k\to 0$ as $k\to\infty$.
\end{enumerate}
and set  $W_k:=(a_k,b_k)\cap \mathbb{Q}$. Then $W_k$ is a clopen subset of $\mathbb{Q}$.

For every $k\in\w$, choose two sequences $\{a_{n,k}\}_{n\in\w}$ and  $\{b_{n,k}\}_{n\in\w}$ of irrational numbers such that
\begin{enumerate}
\item[(c)] $a_k<b_{n+1,k}<a_{n,k}<b_{n,k}< b_k$ for every $n\in\w$;
\item[(d)] $b_{n,k}\to a_k$ as $n\to\infty$,
\end{enumerate}
and set $W_{n,k}:=(a_{n,k},b_{n,k})\cap \mathbb{Q}$. Then for every $n,k\in\w$, $W_{n,k}$ and $\bigcup_{n\in\w} W_{n,k}$ are clopen subsets of $\mathbb{Q}$. It is easy to see that if $0\not=x\in\mathbb{Q}$, then $x$ has a neighborhood which intersects with at most one of the sets $W_{i,j}$, $i,j\in\w$.

For every $n,k\in\w$, we set
\[
U_{n,k}:=\{x_{n,k}\}\times W_{n,k} \;\; \mbox{ and }\;\;  U:=\bigcup_{n,k\in\w} U_{n,k}.
\]
It is clear that all $U_{n,k}$ are clopen subsets of $X$, and $U$ is an open subset of $X$.
\smallskip

{\em Claim 1. $\overline{U}\SM U=\{(x_\infty,0)\}$.}
\smallskip

To prove the claim, let $(y,z)\in \overline{U}\SM U$. If $y=x_{n,k}$ for some $n,k\in\w$, then $y=x_{n,k}$ is isolated in $V(\w)$. Therefore $(\{x_{n,k}\}\times \mathbb{Q})\cap U= \{x_{n,k}\}\times W_{n,k}$ is a closed subset of $U$. Hence $(y,x)\in U$, a contradiction. Thus $y=x_\infty$.

Assume that $z\not= 0$. Choose a neighborhood $\mathcal{O}$ of $z$ which intersects with at most one of the sets $W_{i,j}$, say  $W_{n,k}$. Choose a neighborhood $V$ of $x_\infty$ such that $x_{n,k}\not\in V$. Then $(V\times \mathcal{O})\cap U=\emptyset$, and hence $(x_\infty,z)\not\in \overline{U}$, a contradiction. Thus $z=0$.

It remains to show that $(x_\infty,0)\in \overline{U}\SM U$. Note first that, by construction,  $(x_\infty,0)\not\in U$. We prove that  $(x_\infty,z)\in \overline{U}$. Let $P$ be a neighborhood of $x_\infty$ in $V(\w)$, and $R=(-r,r)\cap \mathbb{Q}$ be a neighborhood of $0$ in $\mathbb{Q}$. Choose $k\in\w$  such that $b_k<r$, so that $W_{n,k}\subseteq R$ for every $n\in\w$. For the chosen $k$, take $n\in\w$ such that $x_{n,k}\in P$. Then $U_{n,k}\subseteq P\times R$. Thus  $(x_\infty,0)\in \overline{U}$, and Claim 1 is proved.
\smallskip

To finish the proof, by Claim 1, it suffices to show that there is no sequences in $U$ converging to $(x_\infty,0)$. Suppose for a contradiction that there is a sequence $\big\{ \big(x_{n_i,k_i},z_{n_i,k_i}\big)\big\}_{i\in\w}$ in $U$ which converges to  $(x_\infty,0)$. Then $x_{n_i,k_i}\to x_\infty$. Therefore there is $k\in\w$ such that $k_i\leq k$ for every $i\in\w$. So, without loss of generality we can assume that $k_i=k$ and $x_{n_i,k}\to x_\infty$. This implies that $z_{n_i,k}\in W_{n_i,k} \subseteq W_k$, and hence  $z_{n_i,k}\not\to 0$ as $i\to\infty$. Therefore $\big(x_{n_i,k_i},z_{n_i,k_i}\big)\not\to (x_\infty,0)$, a contradiction. \qed
\end{proof}

\end{document}
